% ============================================================
%  CHAPTER 15: HEAT ENERGY AND CALORIMETRY
%  The Physics Codex: A Complete University Foundation
%  Korede Samuel Omotosho (Falcon)
%  Aurikrex Learning Systems, 2026
%
%  File: chapters/part2/ch15_calorimetry.tex
% ============================================================

\chapter{Heat Energy and Calorimetry}
\label{ch:calorimetry}
\minitoc
\newpage

\chapterquote{The moving power of mathematical invention
is not reasoning but imagination.}{Augustus de Morgan}

\begin{objectivebox}
By the end of this chapter, you should be able to:
\begin{enumerate}[noitemsep]
  \item Distinguish clearly between heat and
  temperature
  \item Define specific heat capacity and apply it
  to calculate heat transfer
  \item Describe and apply the method of mixtures
  and electrical heating method
  \item Define specific latent heat of fusion and
  vaporisation
  \item Interpret and construct heating and cooling
  curves
  \item Explain why temperature remains constant
  during phase changes using a molecular model
  \item Solve calorimetry problems combining
  sensible heat and latent heat
  \item Explain why coastal regions have more
  moderate climates than inland regions
\end{enumerate}
\end{objectivebox}

% ============================================================
\section{Intuition: Heat and Temperature Are Not the
Same Thing}
% ============================================================

Touch a metal spoon left in a hot pot and a wooden
spoon in the same pot. The metal spoon feels much
hotter. But both are at the same temperature. The
difference is how rapidly each conducts heat into
your fingers.

Now consider a more striking example. Drop a lit
match into a swimming pool. The match is far hotter
than the pool --- perhaps $800°\text{C}$ vs $25°\text{C}$.
Yet the match cannot raise the pool temperature
noticeably. The reason: the pool contains vastly
more thermal energy than the match, despite being
cooler. Temperature tells you how hot something is.
Thermal energy tells you how much heat energy it
contains. A large cool object can contain more
thermal energy than a small hot object.

This distinction is fundamental to everything in
this chapter.

\begin{definitionbox}[Heat and Internal Energy]
\textbf{Internal energy} $U$ is the total kinetic
and potential energy of all the molecules in a
substance. It depends on both the temperature
\textit{and} the quantity of material.

\textbf{Heat} $Q$ is energy in transfer from a
hotter body to a cooler body due to a temperature
difference. Heat is a process, not a property ---
a body does not ``contain heat''; it contains
internal energy. When we say a body ``gains heat,''
we mean energy is transferred into it.

SI unit of heat: joule (J).
\end{definitionbox}

% ============================================================
\section{Specific Heat Capacity}
% ============================================================

Different substances require different amounts of
energy to raise their temperature by the same amount.
Water is famously resistant to temperature change ---
it takes $4200\ \text{J}$ to raise $1\ \text{kg}$ of
water by $1°\text{C}$. The same energy raises $1\ \text{kg}$
of copper by about $24°\text{C}$.

\begin{definitionbox}[Specific Heat Capacity]
The \textbf{specific heat capacity} $c$ of a substance
is the heat energy required to raise the temperature
of $1\ \text{kg}$ of the substance by $1\ \text{K}$
(or $1°\text{C}$):
\[
  Q = mc\Delta\theta
\]
where:
\begin{itemize}[noitemsep]
  \item $Q$ = heat energy transferred (J)
  \item $m$ = mass of substance (kg)
  \item $c$ = specific heat capacity
  (J\,kg$^{-1}$\,K$^{-1}$)
  \item $\Delta\theta$ = temperature change (K or °C)
\end{itemize}
SI unit: J\,kg$^{-1}$\,K$^{-1}$.
\end{definitionbox}

\begin{table}[H]
\centering
\caption{Specific Heat Capacities of Common Substances}
\label{tab:shc}
\begin{tabular}{ll}
\toprule
\textbf{Substance} &
\textbf{Specific Heat Capacity (J\,kg$^{-1}$\,K$^{-1}$)}\\
\midrule
Water (liquid) & 4200\\
Ice & 2100\\
Steam & 2010\\
Seawater & 3900\\
Ethanol & 2400\\
Aluminium & 900\\
Glass & 840\\
Iron/Steel & 450\\
Copper & 390\\
Lead & 128\\
Mercury & 140\\
Air (at constant pressure) & 1010\\
\bottomrule
\end{tabular}
\end{table}

\begin{realworldbox}[Why Coastal Cities Have Milder Climates]
Water has one of the highest specific heat capacities
of any substance. The Atlantic Ocean adjacent to the
West African coastline absorbs enormous amounts of
solar energy during the day without rising much in
temperature, and releases that energy slowly at night.
This moderates the temperature of coastal cities
like Lagos, Accra, and Abidjan --- they rarely experience
the extreme heat of inland cities like Kano or Maiduguri,
even though they receive similar solar radiation.
Kano in the Sahel region, far from the ocean and
surrounded by dry land (low heat capacity), swings
between much higher daytime temperatures and cooler
nights. The ocean's high specific heat capacity is
the reason coastal climates are milder.
\end{realworldbox}

\begin{examplebox}[15.1]
\stepcounter{workedexample}
Calculate the heat required to raise $3.0\ \text{kg}$
of water from $20°\text{C}$ to $100°\text{C}$.
($c_{water} = 4200\ \text{J\,kg}^{-1}\text{K}^{-1}$)
\end{examplebox}

\begin{solutionbox}
\[
  Q = mc\Delta\theta
  = 3.0 \times 4200 \times (100 - 20)
  = 3.0 \times 4200 \times 80
  = 1\,008\,000\ \text{J}
  = 1008\ \text{kJ} \approx 1.0\ \text{MJ}
\]
This is the energy your gas cooker or electric ring
must supply to bring a full pot of water to boiling.
\end{solutionbox}

% ============================================================
\section{Measuring Specific Heat Capacity}
% ============================================================

\subsection{The Electrical Method}

\begin{processbox}[Electrical Method for Specific Heat Capacity]
A heater of known power $P$ (in watts) heats a
mass $m$ of the substance for time $t$.

Heat supplied by heater:
\[
  Q = Pt = VIt
\]
Heat absorbed by substance:
\[
  Q = mc\Delta\theta
\]
Equating (assuming no heat losses):
\[
  mc\Delta\theta = VIt
  \implies c = \frac{VIt}{m\Delta\theta}
\]
\end{processbox}

\begin{diagbox}[Electrical Calorimetry Setup]
\begin{center}
\begin{tikzpicture}[scale=1.0, font=\small]

  % Calorimeter container
  \fill[gray!20] (1.5,0.5) rectangle (6.5,4.5);
  \draw[very thick] (1.5,0.5) rectangle (6.5,4.5);

  % Insulation (lagging)
  \fill[codexgold!20] (1.0,0.2) rectangle (7.0,5.0);
  \fill[gray!20] (1.5,0.5) rectangle (6.5,4.5);
  \draw[thick, dashed, codexgold]
    (1.0,0.2) rectangle (7.0,5.0);
  \node[codexgold, font=\tiny] at (0.5,2.5)
    {Insulation};

  % Liquid/solid being heated
  \fill[codexblue!25] (1.8,0.8) rectangle (6.2,3.2);
  \draw[codexblue] (1.8,3.2) -- (6.2,3.2);
  \node[codexblue, font=\small] at (4.0,2.0)
    {\shortstack{Mass $m$\\(substance)}};

  % Heater coil inside
  \draw[very thick, codexred, decorate,
    decoration={coil, amplitude=4pt, segment length=5pt}]
    (2.5,1.5) -- (5.5,1.5);
  \node[codexred, font=\tiny] at (4.0,1.0)
    {\shortstack{Heater coil\\(power $P = VI$)}};

  % Heater wires
  \draw[thick, codexred] (2.5,1.5) -- (2.5,5.0);
  \draw[thick, codexred] (5.5,1.5) -- (5.5,5.0);

  % Thermometer
  \draw[thick] (4.0,3.2) -- (4.0,5.5);
  \fill[codexred!60] (3.8,3.2) rectangle (4.2,3.6);
  \node[right, font=\tiny] at (4.2,4.5) {Thermometer};

  % Voltmeter (above)
  \draw[thick] (2.5,5.0) -- (2.5,6.2)
    -- (3.8,6.2);
  \draw[thick] (5.5,5.0) -- (5.5,6.2)
    -- (4.2,6.2);
  \fill[white] (3.8,5.8) rectangle (4.2,6.6);
  \draw[thick] (3.8,5.8) rectangle (4.2,6.6);
  \node[font=\tiny] at (4.0,6.2) {V};

  % Ammeter (side)
  \draw[thick] (6.5,2.5) -- (8.0,2.5);
  \draw[thick] (5.5,5.0) -- (6.5,5.0) -- (6.5,2.5);
  \fill[white] (7.6,2.1) rectangle (8.4,2.9);
  \draw[thick] (7.6,2.1) rectangle (8.4,2.9);
  \node[font=\tiny] at (8.0,2.5) {A};

  % Power supply
  \fill[gray!40] (8.0,4.0) rectangle (9.5,5.5);
  \draw[thick] (8.0,4.0) rectangle (9.5,5.5);
  \node[font=\tiny, align=center] at (8.75,4.75)
    {Power\\supply};
  \draw[thick] (8.75,5.5) -- (8.75,6.0)
    -- (4.2,6.0);
  \draw[thick] (8.75,4.0) -- (8.75,3.5)
    -- (8.0,3.5) -- (8.0,2.5);

  % Formula annotation
  \node[align=center, font=\small] at (4.0,-0.3)
    {$c = \dfrac{VIt}{m\Delta\theta}$};

\end{tikzpicture}
\end{center}
\end{diagbox}

\begin{examplebox}[15.2]
\stepcounter{workedexample}
An electrical heater of resistance $20\ \Omega$ is
connected to a $12\ \text{V}$ supply and immersed in
$500\ \text{g}$ of oil. After $3.0\ \text{min}$, the
temperature rises from $22°\text{C}$ to $62°\text{C}$.\\
(a) Calculate the current in the heater.\\
(b) Calculate the heat supplied.\\
(c) Calculate the specific heat capacity of the oil.\\
(d) Why is the experimental value likely to be
higher than the true value?
\end{examplebox}

\begin{solutionbox}
\textbf{(a)} Current:
\[
  I = \frac{V}{R} = \frac{12}{20} = 0.60\ \text{A}
\]

\textbf{(b)} Heat supplied ($t = 3.0 \times 60
= 180\ \text{s}$):
\[
  Q = VIt = 12 \times 0.60 \times 180
  = 1296\ \text{J}
\]

\textbf{(c)} Specific heat capacity:
\[
  c = \frac{Q}{m\Delta\theta}
  = \frac{1296}{0.500 \times 40}
  = \frac{1296}{20}
  = 64.8\ \text{J\,kg}^{-1}\text{K}^{-1}
\]

\textbf{(d)} In a real experiment, some heat is
always lost to the surroundings and to heating the
calorimeter and heater themselves. The oil absorbs
less heat than the $1296\ \text{J}$ supplied,
so the true temperature rise is smaller. Using the
actual temperature rise (which is smaller than it
should be) in the formula makes $c$ appear larger.
To minimise this error: use an insulated container,
heat quickly, and use a correction for heat losses.
\end{solutionbox}

\subsection{The Method of Mixtures}

\begin{processbox}[Method of Mixtures]
A hot solid of known temperature is added to cold
water in an insulated calorimeter. Heat flows from
the hot solid to the cold water until thermal
equilibrium is reached at a common final temperature
$\theta_f$.

Assuming no heat losses:
\[
  \text{Heat lost by hot solid}
  = \text{Heat gained by cold water (+ calorimeter)}
\]
\[
  m_s c_s(\theta_s - \theta_f)
  = m_w c_w(\theta_f - \theta_w)
  + m_c c_c(\theta_f - \theta_w)
\]
where subscripts $s$, $w$, $c$ denote solid,
water, and calorimeter respectively.
\end{processbox}

\begin{diagbox}[Method of Mixtures: Before and After]
\begin{center}
\begin{tikzpicture}[scale=0.95, font=\small]

  % === BEFORE ===
  \begin{scope}[shift={(0,0)}]
    \node[font=\bfseries] at (2.0,5.8) {Before mixing};

    % Hot solid (above)
    \fill[codexred!50] (0.8,4.0) rectangle (3.2,5.0);
    \draw[thick] (0.8,4.0) rectangle (3.2,5.0);
    \node[white, font=\small] at (2.0,4.5)
      {Hot solid};
    \node[codexred, font=\small] at (2.0,3.6)
      {$m_s$, $c_s$, $\theta_s$};
    % Heat waves
    \foreach \x in {1.2,1.8,2.4,3.0} {
      \draw[codexred, thick] (\x,5.0)
        .. controls (\x-0.1,5.2) and (\x+0.1,5.4)
        .. (\x,5.5);
    }

    % Calorimeter with water
    \fill[codexblue!10] (0.3,0.3) rectangle (3.7,3.2);
    \draw[very thick] (0.3,0.3) rectangle (3.7,3.2);
    \fill[codexblue!30] (0.6,0.6) rectangle (3.4,2.5);
    \draw[codexblue] (0.6,2.5) -- (3.4,2.5);
    \node[codexblue, font=\small] at (2.0,1.5)
      {Cold water};
    \node[codexblue, font=\small] at (2.0,0.9)
      {$m_w$, $c_w$, $\theta_w$};
    \node[font=\tiny] at (2.0,2.9)
      {Calorimeter: $m_c$, $c_c$};
  \end{scope}

  % Arrow
  \draw[->, very thick, codexgold]
    (4.1,2.5) -- (5.0,2.5)
    node[above, font=\small]{\shortstack{drop\\solid in}};

  % === AFTER ===
  \begin{scope}[shift={(5.5,0)}]
    \node[font=\bfseries] at (2.0,5.8) {After mixing};

    % Calorimeter with solid and water
    \fill[codexblue!10] (0.3,0.3) rectangle (3.7,3.2);
    \draw[very thick] (0.3,0.3) rectangle (3.7,3.2);
    \fill[codexblue!25] (0.6,0.6) rectangle (3.4,2.8);
    \draw[codexblue] (0.6,2.8) -- (3.4,2.8);
    % Solid submerged
    \fill[codexred!30] (1.2,0.9) rectangle (2.8,1.8);
    \draw[thick] (1.2,0.9) rectangle (2.8,1.8);
    \node[white, font=\tiny] at (2.0,1.35) {solid};
    % Equilibrium label
    \node[codexgold, font=\small] at (2.0,4.4)
      {Thermal equilibrium};
    \node[codexgold, font=\small] at (2.0,4.0)
      {$\theta_f$ (common temperature)};

    % Energy balance equation
    \node[align=center, font=\small] at (2.0,5.3)
      {Heat lost = Heat gained};
    \node[align=center, font=\tiny, codexblue]
      at (2.0,4.8)
      {$m_sc_s(\theta_s-\theta_f)
      = (m_wc_w+m_cc_c)(\theta_f-\theta_w)$};
  \end{scope}

\end{tikzpicture}
\end{center}
\end{diagbox}

\begin{examplebox}[15.3]
\stepcounter{workedexample}
A $200\ \text{g}$ iron block at $90°\text{C}$ is
placed into $300\ \text{g}$ of water at $20°\text{C}$
in an aluminium calorimeter of mass $100\ \text{g}$.
Find the final equilibrium temperature.\\
($c_{iron} = 450$, $c_{water} = 4200$,
$c_{Al} = 900$, all in J\,kg$^{-1}$\,K$^{-1}$)
\end{examplebox}

\begin{solutionbox}
Let final temperature $= \theta_f$.

Heat lost by iron:
$Q_{lost} = 0.200 \times 450 \times (90 - \theta_f)
= 90(90 - \theta_f)$

Heat gained by water:
$Q_{water} = 0.300 \times 4200 \times (\theta_f - 20)
= 1260(\theta_f - 20)$

Heat gained by calorimeter:
$Q_{cal} = 0.100 \times 900 \times (\theta_f - 20)
= 90(\theta_f - 20)$

Setting heat lost = heat gained:
\[
  90(90 - \theta_f) = 1260(\theta_f - 20)
  + 90(\theta_f - 20)
\]
\[
  8100 - 90\theta_f = 1350(\theta_f - 20)
\]
\[
  8100 - 90\theta_f = 1350\theta_f - 27\,000
\]
\[
  8100 + 27\,000 = 1350\theta_f + 90\theta_f
\]
\[
  35\,100 = 1440\theta_f
\]
\[
  \theta_f = \frac{35\,100}{1440} = 24.4°\text{C}
\]

The small iron block at $90°\text{C}$ only raises the
water temperature from $20°\text{C}$ to $24.4°\text{C}$
--- because water has such a high specific heat capacity
that it absorbs a lot of heat for a small temperature rise.
\end{solutionbox}

% ============================================================
\section{Latent Heat}
% ============================================================

Heat the ice cube in a glass of iced water. The
temperature stays at $0°\text{C}$ until all the ice
has melted. Only then does the liquid water begin to
warm up. The heat added during melting seems to
``disappear'' --- it does not raise the temperature.
Where does it go?

It goes into breaking the bonds between molecules.
When ice melts, the rigid crystalline structure of
ice (where each water molecule is held firmly in
place by hydrogen bonds) breaks down. Molecules gain
the freedom to move past each other. This requires
energy --- not to move faster (which would raise
temperature) but to break free from their neighbours
(which changes the state without changing the
temperature).

This energy is called \textbf{latent heat}.

\begin{definitionbox}[Specific Latent Heat]
The \textbf{specific latent heat} $L$ of a substance
is the heat energy required to change the state of
$1\ \text{kg}$ of the substance at constant
temperature:
\[
  Q = mL
\]
SI unit: J\,kg$^{-1}$.

\textbf{Specific latent heat of fusion} $L_f$:
heat required to melt (or freeze) $1\ \text{kg}$.

\textbf{Specific latent heat of vaporisation} $L_v$:
heat required to vaporise (or condense) $1\ \text{kg}$.

For water:
$L_f = 3.36\times10^5\ \text{J\,kg}^{-1}$
(at $0°\text{C}$),
$L_v = 2.26\times10^6\ \text{J\,kg}^{-1}$
(at $100°\text{C}$).
\end{definitionbox}

\begin{keybox}[Why $L_v \gg L_f$]
The latent heat of vaporisation is much larger than
the latent heat of fusion because:
\begin{itemize}[noitemsep]
  \item Melting only breaks some intermolecular bonds
  (molecules remain close together in the liquid)
  \item Vaporisation must completely overcome
  \textit{all} intermolecular attractions (molecules
  must separate to large distances in the gas)
  \item Vaporisation also requires work done against
  atmospheric pressure as the gas expands dramatically
  ($1\ \text{kg}$ of water at $100°\text{C}$
  becomes $1.67\ \text{m}^3$ of steam at atmospheric
  pressure)
\end{itemize}
For water, $L_v \approx 6.7 \times L_f$.
\end{keybox}

% ============================================================
\section{The Heating Curve}
% ============================================================

\begin{diagbox}[Complete Heating Curve: Ice to Steam]
\begin{center}
\begin{tikzpicture}
\begin{axis}[width=13cm, height=7.5cm,
  xlabel={Heat energy supplied $Q$ (kJ)},
  ylabel={Temperature $\theta$ (\textdegree C)},
  xmin=0, xmax=3350,
  ymin=-30, ymax=130,
  grid=major, grid style={gray!15},
  xtick={0,210,630,2890,3140},
  xticklabel style={font=\tiny, text width=0.7cm, align=center},
  xticklabels={0,210,630,2890,3140},
  ytick={-20,0,20,40,60,80,100,120},
  clip=false]
  % Phase 1: Heating ice
  \addplot[codexblue, very thick, domain=0:210]{-20 + (20/210)*x};
  % Phase 2: Melting plateau
  \addplot[codexblue!70, very thick, domain=210:630]{0};
  % Phase 3: Heating water
  \addplot[codexred, very thick, domain=630:2890]{(100/2260)*(x-630)};
  % Phase 4: Vaporising plateau
  \addplot[codexred!70, very thick, domain=2890:3140]{100};
  % Phase 5: Heating steam
  \addplot[codexgold, very thick, domain=3140:3350]{100 + (20/210)*(x-3140)};

  % === LEFT ZONE: compact single-line label in the open wedge, melting label raised clear ===
  \node[codexblue, font=\tiny, align=left, anchor=west] at (axis cs:125,-11)
    {Heating ice\\
    $Q=mc_{ice}\Delta\theta$};
  \node[codexblue!80, font=\small, align=center] at (axis cs:420,55)
    {Melting\\(latent heat)\\$Q=mL_f$};
  \draw[->, codexblue!50, thin] (axis cs:420,40) -- (axis cs:420,3);
  \draw[<->, codexblue!70, thick] (axis cs:210,12) -- (axis cs:630,12);
  \node[below, codexblue, font=\tiny] at (axis cs:420,12){$mL_f = 336$ kJ};

  % === CENTRE ===
  \node[codexred, font=\small, align=center] at (axis cs:1760,50)
    {Heating liquid water\\$Q=mc_w\Delta\theta$};

  % === RIGHT ZONE: each element on its own row ===
  \node[codexred!80, font=\tiny, align=left] at (axis cs:2955,52)
    {Vaporising\\$Q=mL_v$};
  \draw[->, codexred!70, thick] (axis cs:3005,62) -- (axis cs:3005,99);
  \draw[<->, codexred!70, thick] (axis cs:2880,93) -- (axis cs:3130,93);
  \node[below, codexred, font=\tiny] at (axis cs:3015,91){$mL_v = 2260$ kJ};
  \node[codexgold, font=\tiny, anchor=west] at (axis cs:2750,118)
    {Vapour (steam)};

  % Temperature reference lines
  \draw[dashed, gray] (axis cs:0,0) -- (axis cs:210,0);
  \draw[dashed, gray] (axis cs:0,100) -- (axis cs:2890,100);
\end{axis}
\end{tikzpicture}
\end{center}
\small Heating 1 kg of water from $-20$\textdegree C
(ice) to above $100$\textdegree C (steam). The two flat
plateaux are where energy goes into breaking molecular
bonds (latent heat) without raising temperature.
Note how much longer the vaporisation plateau is ---
$L_v$ is over six times $L_f$ for water.
\end{diagbox}

\begin{examplebox}[15.4]
\stepcounter{workedexample}
Calculate the total heat energy required to convert
$500\ \text{g}$ of ice at $-10°\text{C}$ completely
into steam at $100°\text{C}$.\\
($c_{ice} = 2100$, $c_{water} = 4200$
J\,kg$^{-1}$\,K$^{-1}$;
$L_f = 3.36\times10^5$,
$L_v = 2.26\times10^6$ J\,kg$^{-1}$)
\end{examplebox}

\begin{solutionbox}
$m = 0.500\ \text{kg}$. The process has four stages:

\textbf{Stage 1:} Heating ice from $-10°\text{C}$
to $0°\text{C}$:
\[
  Q_1 = mc_{ice}\Delta\theta
  = 0.500 \times 2100 \times 10
  = 10\,500\ \text{J}
\]

\textbf{Stage 2:} Melting ice at $0°\text{C}$:
\[
  Q_2 = mL_f
  = 0.500 \times 3.36\times10^5
  = 168\,000\ \text{J}
\]

\textbf{Stage 3:} Heating water from $0°\text{C}$
to $100°\text{C}$:
\[
  Q_3 = mc_{water}\Delta\theta
  = 0.500 \times 4200 \times 100
  = 210\,000\ \text{J}
\]

\textbf{Stage 4:} Vaporising water at $100°\text{C}$:
\[
  Q_4 = mL_v
  = 0.500 \times 2.26\times10^6
  = 1\,130\,000\ \text{J}
\]

\textbf{Total:}
\[
  Q_{total} = 10\,500 + 168\,000 + 210\,000
  + 1\,130\,000
  = 1\,518\,500\ \text{J} \approx 1.52\ \text{MJ}
\]

Note how the vaporisation step alone ($1.13\ \text{MJ}$)
accounts for $74\%$ of the total energy. This is why
steam burns are so much more severe than hot water burns
at the same temperature --- steam releases an enormous
amount of energy when it condenses on your skin.
\end{solutionbox}

\begin{examplebox}[15.5]
\stepcounter{workedexample}
$200\ \text{g}$ of ice at $0°\text{C}$ is added to
$500\ \text{g}$ of water at $60°\text{C}$ in an
insulated container. Find the final temperature.\\
($c_{water} = 4200\ \text{J\,kg}^{-1}\text{K}^{-1}$,
$L_f = 3.36\times10^5\ \text{J\,kg}^{-1}$)
\end{examplebox}

\begin{solutionbox}
\textbf{First, check if all the ice melts.}

Heat available from warm water cooling to $0°\text{C}$:
\[
  Q_{available} = m_w c_w \Delta\theta
  = 0.500 \times 4200 \times 60
  = 126\,000\ \text{J}
\]

Heat needed to melt all the ice:
\[
  Q_{melt} = m_{ice} L_f
  = 0.200 \times 3.36\times10^5
  = 67\,200\ \text{J}
\]

Since $Q_{available} = 126\,000\ \text{J}
> Q_{melt} = 67\,200\ \text{J}$, \textbf{all the
ice melts}.

Remaining heat after melting ice:
\[
  Q_{remaining} = 126\,000 - 67\,200
  = 58\,800\ \text{J}
\]

This remaining heat warms the total water mass
($0.200 + 0.500 = 0.700\ \text{kg}$) from $0°\text{C}$:
\[
  \theta_f = \frac{Q_{remaining}}{m_{total}c_w}
  = \frac{58\,800}{0.700 \times 4200}
  = \frac{58\,800}{2940}
  = 20.0°\text{C}
\]

\textbf{Final temperature: $20.0°\text{C}$}
\end{solutionbox}

% ============================================================
\section{Evaporation vs Boiling}
% ============================================================

\begin{defbox}[Evaporation and Boiling Compared]
\begin{center}
\begin{tabular}{ll}
\toprule
\textbf{Evaporation} & \textbf{Boiling}\\
\midrule
Occurs at \textit{any} temperature & Only at boiling point\\
Occurs at the surface only & Occurs throughout the liquid\\
No bubbles form & Bubbles form throughout\\
Slow process & Rapid process\\
Causes cooling of the liquid & Requires continuous heat supply\\
Rate increases with temperature, & Depends on external pressure\\
surface area, and air flow & \\
\bottomrule
\end{tabular}
\end{center}
\vspace{0.3cm}
\textbf{Why evaporation causes cooling:}
The fastest molecules escape the surface. The average
speed of the remaining molecules decreases. Lower
average speed means lower average KE means lower
temperature. Sweating cools the body by this mechanism.
\end{defbox}

% ============================================================
\section{Chapter Summary}
% ============================================================

\begin{summarybox}
\textbf{Heat $\neq$ Temperature:} heat is energy in
transfer; temperature measures average molecular KE.

\textbf{Specific heat capacity:}
$Q = mc\Delta\theta$
(J\,kg$^{-1}$\,K$^{-1}$)

$c_{water} = 4200\ \text{J\,kg}^{-1}\text{K}^{-1}$
(very high --- moderates coastal climates)

\textbf{Electrical method:}
$c = VIt/(m\Delta\theta)$

\textbf{Method of mixtures:}
heat lost $=$ heat gained (in insulated container)

\textbf{Specific latent heat:}
$Q = mL$ (J\,kg$^{-1}$)

$L_f(water) = 3.36\times10^5\ \text{J\,kg}^{-1}$

$L_v(water) = 2.26\times10^6\ \text{J\,kg}^{-1}$

$L_v \gg L_f$ because vaporisation must overcome
\textit{all} intermolecular bonds.

\textbf{Heating curve:} two plateaux (melting and
boiling) where $T$ is constant despite heat input.

\textbf{Evaporation} cools a liquid (fastest
molecules escape); \textbf{boiling} requires
external heat input.
\end{summarybox}

% ============================================================
\newpage
\begin{exercisebox}[15]

\subsection*{Section A: Multiple Choice}

\textbf{A1.} The heat required to raise $2.0\ \text{kg}$
of copper ($c = 390\ \text{J\,kg}^{-1}\text{K}^{-1}$)
from $20°\text{C}$ to $70°\text{C}$ is:

\mcoption{A}{$3900\ \text{J}$}
\mcoption{B}{$19\,500\ \text{J}$}
\mcoption{C}{$39\,000\ \text{J}$}
\mcoption{D}{$78\,000\ \text{J}$}
\ansref{Ch15-A1}

\vspace{0.4cm}

\textbf{A2.} During melting, the temperature of a
pure substance:

\mcoption{A}{Rises steadily}
\mcoption{B}{Falls steadily}
\mcoption{C}{Remains constant}
\mcoption{D}{Rises then falls}
\ansref{Ch15-A2}

\vspace{0.4cm}

\textbf{A3.} Which of the following explains why
the latent heat of vaporisation is greater than the
latent heat of fusion for water?

\mcoption{A}{More energy is needed to heat the
water than to melt the ice}
\mcoption{B}{Vaporisation requires complete
separation of all molecules; melting only partially
disrupts the bonds}
\mcoption{C}{Water has a higher specific heat
capacity than ice}
\mcoption{D}{The boiling point is higher than
the melting point}
\ansref{Ch15-A3}

\vspace{0.4cm}

\textbf{A4.} A $500\ \text{W}$ heater heats
$2.0\ \text{kg}$ of a liquid for $4.0\ \text{min}$,
raising its temperature by $30°\text{C}$. The
specific heat capacity of the liquid is:

\mcoption{A}{$500\ \text{J\,kg}^{-1}\text{K}^{-1}$}
\mcoption{B}{$1000\ \text{J\,kg}^{-1}\text{K}^{-1}$}
\mcoption{C}{$2000\ \text{J\,kg}^{-1}\text{K}^{-1}$}
\mcoption{D}{$4000\ \text{J\,kg}^{-1}\text{K}^{-1}$}
\ansref{Ch15-A4}

\vspace{0.4cm}

\textbf{A5.} $100\ \text{g}$ of ice at $0°\text{C}$
is mixed with $100\ \text{g}$ of water at $80°\text{C}$
in an insulated container.
($c_w = 4200$, $L_f = 3.36\times10^5\ \text{J\,kg}^{-1}$)
The final temperature is:

\mcoption{A}{$0°\text{C}$ (not all ice melts)}
\mcoption{B}{$40°\text{C}$}
\mcoption{C}{$6.7°\text{C}$}
\mcoption{D}{$80°\text{C}$}
\ansref{Ch15-A5}

\vspace{0.4cm}

\textbf{A6.} Evaporation differs from boiling
because evaporation:

\mcoption{A}{Only occurs at $100°\text{C}$}
\mcoption{B}{Occurs at the surface at any
temperature and cools the liquid}
\mcoption{C}{Requires more energy than boiling}
\mcoption{D}{Produces bubbles throughout the liquid}
\ansref{Ch15-A6}

\vspace{0.4cm}

\textbf{A7.} A calorimeter experiment gives a
specific heat capacity higher than the true value.
The most likely cause is:

\mcoption{A}{The heater was too powerful}
\mcoption{B}{Heat was lost to the surroundings}
\mcoption{C}{The mass was measured incorrectly}
\mcoption{D}{The thermometer was inaccurate}
\ansref{Ch15-A7}

\vspace{0.4cm}

\textbf{A8.} The heat needed to completely vaporise
$200\ \text{g}$ of water at $100°\text{C}$ is:
($L_v = 2.26\times10^6\ \text{J\,kg}^{-1}$)

\mcoption{A}{$45\,200\ \text{J}$}
\mcoption{B}{$226\,000\ \text{J}$}
\mcoption{C}{$452\,000\ \text{J}$}
\mcoption{D}{$2\,260\,000\ \text{J}$}
\ansref{Ch15-A8}

\vspace{0.4cm}

\textbf{A9.} In the method of mixtures, the
principle used is:

\mcoption{A}{Conservation of momentum}
\mcoption{B}{Conservation of energy (heat lost
= heat gained)}
\mcoption{C}{Newton's law of cooling}
\mcoption{D}{Boyle's Law}
\ansref{Ch15-A9}

\vspace{0.4cm}

\textbf{A10.} A steam burn at $100°\text{C}$ is
more severe than a water burn at $100°\text{C}$
because:

\mcoption{A}{Steam is at higher pressure}
\mcoption{B}{Steam moves faster than water}
\mcoption{C}{Steam releases latent heat of
vaporisation when it condenses on the skin}
\mcoption{D}{Steam has a lower specific heat
capacity than water}
\ansref{Ch15-A10}

\vspace{0.6cm}

\subsection*{Section B: Theory and Calculations}

\textbf{B1.} An immersion heater rated at $2.0\ \text{kW}$
is used to heat $5.0\ \text{kg}$ of water from
$15°\text{C}$.
($c_{water} = 4200\ \text{J\,kg}^{-1}\text{K}^{-1}$,
$L_v = 2.26\times10^6\ \text{J\,kg}^{-1}$)\\[4pt]
(a) How long does it take to bring the water to
$100°\text{C}$? Assume no heat losses.\\
(b) The water is left heating. How much additional
time is needed to vaporise $1.0\ \text{kg}$ of water?\\
(c) In practice, the heater takes $12\%$ longer
than calculated. What is the actual efficiency
of the heater-water system?\\
(d) Explain two reasons why the actual time is
longer than the theoretical time.
\hfill\ansref{Ch15-B1}

\vspace{0.4cm}

\textbf{B2.} A $150\ \text{g}$ aluminium calorimeter
contains $300\ \text{g}$ of water at $18°\text{C}$.
A $200\ \text{g}$ steel block is heated to $120°\text{C}$
in boiling water and quickly transferred to the
calorimeter. The final temperature is $28°\text{C}$.\\
($c_{Al} = 900$, $c_{steel} = 450$,
$c_{water} = 4200$, all J\,kg$^{-1}$\,K$^{-1}$)\\[4pt]
(a) Calculate the heat gained by the water and
calorimeter.\\
(b) Calculate the heat lost by the steel block
using your answer to (a).\\
(c) Using the value from (b), calculate the
specific heat capacity of steel.\\
(d) Compare your answer with the accepted value of
$450\ \text{J\,kg}^{-1}\text{K}^{-1}$ and suggest
reasons for any discrepancy.
\hfill\ansref{Ch15-B2}

\vspace{0.4cm}

\textbf{B3.} $400\ \text{g}$ of ice at $-15°\text{C}$
is added to $600\ \text{g}$ of water at $50°\text{C}$
in an insulated container.\\
($c_{ice} = 2100$, $c_{water} = 4200$
J\,kg$^{-1}$\,K$^{-1}$;
$L_f = 3.36\times10^5\ \text{J\,kg}^{-1}$)\\[4pt]
(a) Calculate the heat needed to bring the ice to
$0°\text{C}$.\\
(b) Calculate the heat available from the warm
water cooling to $0°\text{C}$.\\
(c) Determine whether all the ice melts. Show your
reasoning.\\
(d) Calculate the final temperature of the mixture.
\hfill\ansref{Ch15-B3}

\vspace{0.4cm}

\textbf{B4.} A refrigerator removes heat from a
tray of water to make ice cubes.\\
($c_{water} = 4200$, $c_{ice} = 2100$
J\,kg$^{-1}$\,K$^{-1}$;
$L_f = 3.36\times10^5\ \text{J\,kg}^{-1}$)\\[4pt]
(a) Calculate the total heat that must be removed
to convert $500\ \text{g}$ of water at $25°\text{C}$
into ice at $-10°\text{C}$.\\
(b) The refrigerator removes heat at an average
rate of $80\ \text{W}$. How long does it take to
produce the ice?\\
(c) Sketch the cooling curve for this process,
labelling all stages and temperatures.\\
(d) Why does the freezing process take longer than
simply cooling the same mass of water from
$25°\text{C}$ to $-10°\text{C}$ without any
phase change?
\hfill\ansref{Ch15-B4}

\vspace{0.4cm}

\textbf{B5.} (a) Explain at the molecular level
why:
(i) evaporation causes cooling of the liquid;
(ii) the boiling point of water decreases at
higher altitudes;
(iii) pressure-cooking at higher-than-atmospheric
pressure raises the boiling point of water.\\[4pt]
(b) In a pressure cooker operating at
$1.5\times10^5\ \text{Pa}$ (above atmospheric),
the effective boiling point of water is approximately
$111°\text{C}$.
Calculate the heat needed to convert $300\ \text{g}$
of water at $20°\text{C}$ to steam at $111°\text{C}$
inside the cooker. (Use the same $L_v$ value as
for $100°\text{C}$ as an approximation.)\\[4pt]
(c) A $70\ \text{kg}$ athlete sweats
$1.5\ \text{L}$ of water during exercise.
Assuming all the water evaporates from the skin
($L_v = 2.4\times10^6\ \text{J\,kg}^{-1}$
at body temperature), calculate the heat removed
from the body. If the athlete's body has specific
heat capacity $3500\ \text{J\,kg}^{-1}\text{K}^{-1}$,
by how much would the body temperature have risen
if sweating had not occurred?
\hfill\ansref{Ch15-B5}

\end{exercisebox}